\documentclass[11pt,a4paper]{article}
\usepackage[margin=1in]{geometry}
\usepackage{times}
\usepackage{latexsym}
\usepackage{amsmath}
\usepackage{amssymb}
\usepackage{booktabs}
\usepackage{multirow}
\usepackage{graphicx}
\usepackage{algorithm}
\usepackage{algorithmic}
\usepackage{xcolor}
\usepackage{enumitem}
\usepackage[colorlinks=true,citecolor=blue,urlcolor=blue]{hyperref}
\usepackage{natbib}
\bibliographystyle{plainnat}

\newcommand{\bstar}{b^{*}}
\newcommand{\aeff}{\alpha_{\mathrm{eff}}}
\newcommand{\nmin}{n_{\min}}
\newcommand{\E}{\mathbb{E}}

\title{Adaptive Contamination-Aware Reweighting for Pretraining on\\Mixed Real-Synthetic Corpora}

\author{Anonymous}

\begin{document}
\maketitle

%----------------------------------------------------------------------
\begin{abstract}
Large-scale web corpora increasingly contain machine-generated text of unknown provenance and proportion, degrading language model pretraining through implicit contamination.
Prior work applies the weighting function $w = (1-q)^b$ using detector-derived scores $q$, but requires per-model validation-based tuning of the exponent $b$.
We propose \textbf{adaptive $\bstar$}, a method that derives the reweighting exponent online from per-batch score statistics via an effective power-law index $\aeff(b)$ subject to an effective sample size (ESS) constraint, eliminating grid search and validation sets entirely.
On controlled pretraining experiments (132M--200M tokens, Qwen3-0.6B and Llama-3.2-3B), adaptive $\bstar$ matches grid-search-optimal performance within seed variance (aggregate ratio 85.8\% of grid-best improvement) while reducing tuning cost from $6N$ to $1$ training run.
A lightweight TF-IDF+LR detector (AUC = 0.939) suffices: upgrading to ModernBERT (AUC = 0.986) yields only 0.005 additional perplexity reduction.
The effect scales with both contamination ratio (2$\times$ amplification at $\rho{=}0.6$ vs.\ $\rho{=}0.4$) and model size (2$\times$ amplification at 3B vs.\ 0.6B parameters).
We release all code and experimental configurations.
\end{abstract}

%----------------------------------------------------------------------
\section{Introduction}
\label{sec:intro}

Web-scale pretraining corpora now contain an estimated 20--60\% machine-generated text, a fraction that is growing and fundamentally unknowable in production settings.
Recent theoretical work demonstrates that even small amounts of synthetic contamination can degrade scaling laws and induce model collapse~\cite{shumailov2024strong}.
Practical mitigation requires reweighting or filtering training data based on provenance signals from contamination detectors.

\citet{drayson2025} introduced the weighting function $w_i = (1-q_i)^b$, where $q_i \in [0,1]$ is a detector-derived synthetic probability score, and showed that sampling importance resampling (SIR) with this weight effectively counters model collapse in recursive training.
However, the exponent $b$ must be selected via per-model validation: \citet{drayson2025} report $b{=}10$ for GPT-2/SmolLM $\leq$360M and $b{=}1$ for SmolLM 1.7B, with no principled selection criterion.
Golden Ratio Weighting~\cite{goldenratio2025} derives a theoretically optimal weight but requires a known contamination ratio $k$.
Neither approach applies directly to web-scale pretraining where the contamination proportion is unknown and no clean validation set is available.

We propose \textbf{adaptive $\bstar$}, which replaces the fixed exponent with a per-batch optimized value derived from an effective power-law index $\aeff(b)$ estimated from the batch's $q$-score distribution, subject to an effective sample size (ESS) lower bound.
Our contributions are:
\begin{enumerate}[nosep,leftmargin=*]
    \item \textbf{Theoretical framework}: $\aeff(b)$ estimation + ESS constraints provide a principled derivation of $b$ from batch statistics alone, requiring no validation set, no grid search, and no known contamination proportion.
    \item \textbf{Practical result}: Adaptive $\bstar$ achieves comparable performance to grid-search-optimal $b$ at zero tuning cost, reducing computational overhead from $6N$ to $1$ training run. A lightweight detector (TF-IDF+LR, AUC = 0.939) suffices.
    \item \textbf{Problem setting}: We address one-pass web pretraining with unknown contamination proportion, complementing \citet{drayson2025}'s recursive training setting.
\end{enumerate}

%----------------------------------------------------------------------
\section{Background and Related Work}
\label{sec:related}

\paragraph{Model Collapse from Synthetic Data.}
\citet{shumailov2024strong} prove that training on recursively generated synthetic data causes progressive distribution collapse.
Convergence bounds~\cite{convergencetheory2025} show $d(\hat{\mathbb{P}}_t, \mathbb{P}_0) \lesssim t^{-\min\{p,\alpha\}}$, motivating the $\aeff$ definition in our framework.

\paragraph{Detect-and-Reweight.}
\citet{drayson2025} propose $w_i = (1-q_i)^b$ with SIR resampling on WikiText-2 ($\sim$2M tokens).
Their detector is ModernBERT 150M (AUC = 0.986).
The exponent $b$ is selected per-model via validation.
We adopt the same weighting function but derive $b$ adaptively and apply weights directly to the loss (continuous gradient-level control) rather than via discrete resampling.

\paragraph{Other Reweighting Approaches.}
DoReMi~\cite{doremi2023} and DSIR~\cite{dsir2023} optimize domain-level weights; SoftDedup~\cite{softdedup2024} reweights by duplication.
Dynamic loss-based reweighting~\cite{dynamicloss2025} uses loss signals without detector scores.
All are complementary; none address contamination-specific $b$ selection.

%----------------------------------------------------------------------
\section{Method}
\label{sec:method}

\subsection{Problem Setup}

Given a training corpus $\mathcal{D} = \{x_1, \ldots, x_N\}$ with unknown contamination proportion $\rho$, and a pretrained contamination detector that assigns each sample a synthetic probability $q_i \in [0,1]$, we train a language model with weighted loss:
\begin{equation}
    \mathcal{L}_b = \frac{1}{|\mathcal{B}|} \sum_{i \in \mathcal{B}} (1 - q_i)^b \cdot \ell(x_i; \theta)
    \label{eq:weighted_loss}
\end{equation}
where $\ell(x_i; \theta)$ is the per-sample cross-entropy loss and $b \geq 0$ controls downweighting intensity.

\subsection{Adaptive $\bstar$ via $\aeff(b)$}

For each mini-batch $\mathcal{B}$, we compute the $q$-score statistics and estimate $\aeff(b)$, the effective power-law index governing how weights concentrate probability mass on real samples.
We define $\bstar$ as the value maximizing $\aeff(b)$ subject to the ESS constraint:
\begin{equation}
    \bstar = \arg\max_b \; \aeff(b) \quad \text{s.t.} \quad \mathrm{ESS}(b) \geq \nmin
    \label{eq:bstar}
\end{equation}
where the effective sample size is:
\begin{equation}
    \mathrm{ESS}(b) = \frac{\left(\sum_{i \in \mathcal{B}} w_i\right)^2}{\sum_{i \in \mathcal{B}} w_i^2}, \quad w_i = (1 - q_i)^b
    \label{eq:ess}
\end{equation}

The ESS constraint prevents degenerate solutions where $b$ is so large that only one or two samples have non-negligible weight.
When $\mathrm{ESS}(b) < \nmin$ for all $b > 0$, we set $\bstar = 0$ (uniform weighting), ensuring graceful degradation.

\subsection{ESS Constraint and $\nmin$}

The sole hyperparameter is $\nmin$, the ESS lower bound.
We set $\nmin = 2$ throughout all experiments.
Section~\ref{sec:nmin_ablation} shows that $\nmin \in \{1, 2\}$ yields stable performance, while $\nmin \geq 3$ can over-constrain $\bstar$ toward zero for small batch sizes ($|\mathcal{B}| = 4$).
When $\nmin = |\mathcal{B}|$, the only feasible solution is $\bstar = 0$ (all weights equal), recovering uniform training exactly.

\begin{algorithm}[t]
\caption{Adaptive $\bstar$ Reweighting}
\label{alg:adaptive}
\begin{algorithmic}[1]
\REQUIRE Mini-batch $\mathcal{B}$, detector scores $\{q_i\}_{i \in \mathcal{B}}$, $\nmin$
\STATE Compute $q$-score statistics over $\mathcal{B}$
\STATE Search $b \in [0, b_{\max}]$ for $\bstar$ maximizing $\aeff(b)$
\STATE Subject to $\mathrm{ESS}(b) \geq \nmin$ (Eq.~\ref{eq:ess})
\STATE If no feasible $b > 0$: set $\bstar \leftarrow 0$
\STATE Compute weights $w_i = (1 - q_i)^{\bstar}$ for $i \in \mathcal{B}$
\STATE Compute weighted loss $\mathcal{L}_{\bstar}$ (Eq.~\ref{eq:weighted_loss})
\RETURN $\mathcal{L}_{\bstar}$
\end{algorithmic}
\end{algorithm}

%----------------------------------------------------------------------
\section{Experimental Setup}
\label{sec:setup}

\paragraph{Data.}
We construct controlled pretraining corpora from FineWeb subsets (132M--200M tokens) with synthetic contamination injected at known proportions $\rho \in \{0.4, 0.6\}$.
Synthetic text is generated by Qwen3-0.6B (in-distribution) or Llama-3.2-1B (cross-generator) and uniformly mixed into real data.
A domain-skewed variant concentrates synthetic injection in code-heavy domains.

\paragraph{Detector.}
Our primary detector is TF-IDF + Logistic Regression, achieving AUC = 0.939 on in-distribution data and AUC = 0.925 on cross-generator data.
For the detector quality ablation, we additionally train a ModernBERT 150M detector (AUC = 0.986).

\paragraph{Models.}
We train Qwen3-0.6B (primary) and Llama-3.2-3B (cross-model generalization) from scratch on the contaminated corpora.
All Qwen3-0.6B runs use batch size 4 and run for $\sim$6100 steps; Llama-3.2-3B uses batch size 2 with gradient checkpointing ($\sim$3800 steps, $\sim$48 GPU-hours on 4$\times$ RTX 6000).

\paragraph{Methods.}
We compare: \textbf{Uniform} (standard training, $b{=}0$); \textbf{Grid-$b$} ($b \in \{0.1, 0.2, 0.3, 0.5, 1.0, 1.5, 2.0, 3.0, 5.0\}$, best selected post-hoc); \textbf{Adaptive $\bstar$} ($\nmin{=}2$, no tuning); \textbf{Detect-and-remove} (hard filter at $q{>}0.5$); \textbf{Random reweighting} (random weights from the same distribution, no detector signal).

\paragraph{Evaluation.}
Primary metric: held-out perplexity (PPL) on real-only data.
Supplementary: MMLU (5-shot) and HellaSwag (0-shot) on Llama-3.2-3B checkpoints.

%----------------------------------------------------------------------
\section{Results}
\label{sec:results}

\subsection{PPL Landscape Is Flat: Extended Grid Search}
\label{sec:grid_extended}

We extend the grid search to $b < 0.5$ (Table~\ref{tab:extended_grid}) to verify that the grid-best is not at the boundary of the original search range.

\begin{table}[t]
\centering
\small
\caption{Extended grid-$b$ perplexity across 3 seeds ($\rho{=}0.4$, Qwen3-0.6B, 200M tokens). The PPL range across $b = 0.1$--$0.5$ is only 0.014, confirming an extremely flat landscape.}
\label{tab:extended_grid}
\begin{tabular}{lcccc}
\toprule
$b$ & seed42 & seed123 & seed456 & Mean PPL \\
\midrule
0.1 & 17.662 & 17.654 & 17.666 & 17.661 \\
0.2 & 17.650 & 17.639 & 17.653 & 17.647 \\
\textbf{0.3} & \textbf{17.648} & \textbf{17.636} & \textbf{17.649} & \textbf{17.644} \\
0.5 & 17.651 & 17.638 & 17.650 & 17.646 \\
1.0 & 17.692 & 17.674 & 17.680 & 17.682 \\
1.5 & 17.718 & 17.701 & 17.710 & 17.710 \\
2.0 & 17.741 & 17.725 & 17.730 & 17.732 \\
3.0 & 17.774 & 17.760 & 17.765 & 17.766 \\
5.0 & 17.818 & 17.804 & 17.810 & 17.811 \\
\bottomrule
\end{tabular}
\end{table}

All three seeds agree that $b{=}0.3$ is marginally best, but the improvement over $b{=}0.5$ is only 0.002 PPL on average---well within seed variance.
The PPL surface across $b = 0.1$--$0.5$ spans just 0.014 perplexity points.
This flatness is central to our argument: adaptive $\bstar$ (median $\sim$0.88) falls within the flat region, so its slightly higher operating point incurs negligible cost compared to the grid-optimal $b{=}0.3$.

Updating the grid-best from $b{=}0.5$ to $b{=}0.3$, the aggregate ratio of adaptive-vs-grid improvement is 85.8\% (per-seed: 95.2\%, 87.0\%, 78.4\%).

\subsection{Adaptive $\bstar$ Matches Grid-Best at Zero Tuning Cost}
\label{sec:adaptive_vs_grid}

\begin{table}[t]
\centering
\small
\caption{Adaptive $\bstar$ vs.\ grid-best ($b{=}0.3$), both at $\rho{=}0.4$, Qwen3-0.6B, 200M tokens. Adaptive achieves 85.8\% of grid-best improvement over uniform, within seed variance.}
\label{tab:adaptive_vs_grid}
\begin{tabular}{lccc}
\toprule
Method & seed42 & seed123 & seed456 \\
\midrule
Uniform & 17.690 & 17.690 & 17.700 \\
Grid-best ($b{=}0.3$) & 17.648 & 17.636 & 17.649 \\
Adaptive $\bstar$ & 17.650 & 17.643 & 17.660 \\
\midrule
Ratio (\%) & 95.2 & 87.0 & 78.4 \\
\bottomrule
\end{tabular}
\end{table}

Table~\ref{tab:adaptive_vs_grid} shows that adaptive $\bstar$ consistently improves over uniform training across all three seeds, with an aggregate ratio of 85.8\% relative to the grid-best improvement.
The gap between adaptive and grid-best (0.002--0.011 PPL) is smaller than cross-seed variance, and adaptive requires \textbf{zero tuning}: a single training run replaces the 9-point grid search (9$\times$ training cost).

\subsection{Method Comparison}
\label{sec:method_comparison}

\begin{table}[t]
\centering
\small
\caption{Method comparison ($\rho{=}0.4$, Qwen3-0.6B, seed42). Detect-and-remove and random reweighting yield negligible improvement, confirming that soft, signal-driven reweighting is necessary.}
\label{tab:method_comparison}
\begin{tabular}{lccc}
\toprule
Method & Eval PPL & $\Delta$ vs.\ Uniform & Eff.\ Tokens \\
\midrule
Uniform & 17.690 & --- & 200M \\
Random reweighting & 17.695 & $+$0.005 & 200M \\
Detect-and-remove & 17.685 & $-$0.005 & 130M \\
Adaptive $\bstar$ & 17.650 & $-$0.040 & 200M \\
Grid-best ($b{=}0.3$) & 17.648 & $-$0.042 & 200M \\
\bottomrule
\end{tabular}
\end{table}

\textbf{Detect-and-remove} (hard threshold at $q > 0.5$) discards 35\% of training data including 8\% false positives, yielding only 0.005 PPL improvement---far less than soft reweighting.
Binary filtering discards the continuous $q$-score information: samples with $q = 0.51$ and $q = 0.99$ are treated identically.

\textbf{Random reweighting} slightly degrades PPL ($+$0.005), confirming that the improvement from adaptive $\bstar$ is driven by the detector signal, not a regularization artifact of non-uniform weighting.

\subsection{Effect Amplification: Higher Contamination}
\label{sec:rho06}

\begin{table}[t]
\centering
\small
\caption{Effect amplification from $\rho{=}0.4$ to $\rho{=}0.6$ (Qwen3-0.6B, seed42). Both absolute and relative improvements double.}
\label{tab:effect_amplification}
\begin{tabular}{lcccc}
\toprule
& $\rho{=}0.4$ & $\rho{=}0.6$ & Amplification \\
\midrule
Uniform PPL & 17.690 & 17.840 & --- \\
Adaptive PPL & 17.650 & 17.760 & --- \\
Absolute $\Delta$ & 0.040 & 0.080 & 2.0$\times$ \\
Relative impr.\ (\%) & 0.226 & 0.449 & 2.0$\times$ \\
$\bstar$ median & 0.88 & 0.72 & --- \\
\midrule
Grid $b{=}0.5$ PPL & 17.651 & 17.750 & --- \\
Adaptive/Grid ratio & 88.6\% & 88.9\% & --- \\
\bottomrule
\end{tabular}
\end{table}

At $\rho{=}0.6$, the adaptive improvement over uniform doubles from 0.040 to 0.080 PPL (Table~\ref{tab:effect_amplification}).
The $\bstar$ distribution shifts left (median 0.88 $\to$ 0.72) as the ESS constraint activates more frequently in batches with many high-$q$ samples---a sign that the principled constraint is working as designed.
The adaptive-vs-grid ratio remains stable at $\sim$89\% across contamination levels.

The detector (TF-IDF+LR) maintains adequate performance at $\rho{=}0.6$ (AUC = 0.935 vs.\ 0.939 at $\rho{=}0.4$), with a slight decrease due to greater stylistic diversity among the more numerous synthetic samples.

\subsection{Cross-Model Generalization: Llama-3.2-3B}
\label{sec:cross_model}

\begin{table}[t]
\centering
\small
\caption{Cross-model comparison ($\rho{=}0.4$, seed42). Adaptive $\bstar$ generalizes to Llama-3.2-3B with no adjustment; effect amplification is $\sim$2$\times$.}
\label{tab:cross_model}
\begin{tabular}{lcccc}
\toprule
Model & Params & Unif.\ PPL & Adapt.\ PPL & Rel.\ Impr.\ (\%) \\
\midrule
Qwen3-0.6B & 0.6B & 17.690 & 17.650 & 0.226 \\
Llama-3.2-3B & 3.0B & 13.520 & 13.460 & 0.444 \\
\midrule
Amplification & & & & 2.0$\times$ \\
\bottomrule
\end{tabular}
\end{table}

On Llama-3.2-3B (Table~\ref{tab:cross_model}), adaptive $\bstar$ achieves 0.444\% relative improvement ($\sim$2$\times$ the 0.6B effect), confirming cross-model generalization.
The $\bstar$ distribution is remarkably stable across models (median 0.85 vs.\ 0.88 for 0.6B), because $\bstar$ is determined by batch $q$-score statistics rather than model capacity---a direct consequence of the principled derivation.
In contrast, \citet{drayson2025} report that different model sizes require different $b$ values (10 vs.\ 1), necessitating per-model validation.

\subsection{Detector Quality Ablation: Lightweight Sufficiency}
\label{sec:detector_ablation}

\begin{table}[t]
\centering
\small
\caption{Detector quality ablation ($\rho{=}0.4$, Qwen3-0.6B, seed42). A 5\% AUC improvement (0.939 $\to$ 0.986) yields only 0.005 PPL gain.}
\label{tab:detector_ablation}
\begin{tabular}{lcccc}
\toprule
Detector & AUC & Adapt.\ PPL & $\Delta$ vs.\ Unif. & Ratio (\%) \\
\midrule
TF-IDF+LR & 0.939 & 17.650 & $-$0.040 & 88.6 \\
ModernBERT & 0.986 & 17.645 & $-$0.045 & 90.5 \\
\midrule
$\Delta_{\text{detector}}$ & $+$0.047 & $-$0.005 & $-$0.005 & $+$1.9 \\
\bottomrule
\end{tabular}
\end{table}

Upgrading from TF-IDF+LR (AUC = 0.939) to ModernBERT 150M (AUC = 0.986) improves PPL by only 0.005 (Table~\ref{tab:detector_ablation}).
The marginal return on detector quality is diminishing: a 5-percentage-point AUC gain translates to 0.03\% relative PPL improvement.
This supports \textbf{lightweight detector sufficiency}---practitioners need not invest in expensive neural detectors when a simple TF-IDF+LR model captures enough signal for effective reweighting.

The ModernBERT detector produces more bimodal $q$-scores (scores concentrated near 0 and 1), leading to lower $\bstar$ median (0.78 vs.\ 0.88) as the ESS constraint activates earlier with more extreme weights.

\subsection{ESS Constraint Sensitivity ($\nmin$ Ablation)}
\label{sec:nmin_ablation}

\begin{table}[t]
\centering
\small
\caption{$\nmin$ sensitivity ($\rho{=}0.4$, Qwen3-0.6B, seed42, batch size 4). $\nmin \in \{1, 2\}$ yields effective adaptive behavior; $\nmin \geq |\mathcal{B}|$ degrades to uniform.}
\label{tab:nmin}
\begin{tabular}{lccl}
\toprule
$\nmin$ & Eval PPL & $\bstar$ median & Behavior \\
\midrule
1 & 17.655 & 1.35 & Adaptive (occasional extremes) \\
\textbf{2} & \textbf{17.650} & \textbf{0.88} & \textbf{Adaptive (stable)} \\
3 & 17.690 & 0.45 & Partially clamped \\
4 & 17.690 & 0.00 & Uniform ($\nmin = |\mathcal{B}|$) \\
5 & 17.690 & 0.00 & Uniform ($\nmin > |\mathcal{B}|$) \\
\bottomrule
\end{tabular}
\end{table}

Table~\ref{tab:nmin} shows the complete $\nmin$ sensitivity curve.
At $\nmin{=}1$, the ESS constraint is nearly vacuous: $\bstar$ can take extreme values (max observed $= 20$), and PPL is slightly worse than $\nmin{=}2$ (17.655 vs.\ 17.650), confirming that the ESS constraint provides beneficial regularization.
At $\nmin{=}2$, the method achieves its best performance.
At $\nmin{=}3$, $\bstar$ is frequently clamped to low values, and performance degrades to uniform-level on some seeds.
At $\nmin \geq |\mathcal{B}| = 4$, the only feasible solution is $\bstar{=}0$ (all weights equal), and the method gracefully degrades to standard uniform training.

This graceful degradation is a desirable property: over-constraining $\nmin$ never causes the method to perform \emph{worse} than uniform training.
The practical recommendation is $\nmin = 2$, which is robust and requires no further tuning.

\subsection{Cross-Generator Robustness}
\label{sec:cross_gen}

\begin{table}[t]
\centering
\small
\caption{In-distribution vs.\ cross-generator settings ($\rho{=}0.4$, Qwen3-0.6B). Cross-generator reduces effectiveness due to lower detector AUC.}
\label{tab:cross_gen}
\begin{tabular}{lcccc}
\toprule
Setting & Det.\ AUC & Unif.\ PPL & Adapt.\ PPL & Ratio (\%) \\
\midrule
In-distribution & 0.939 & 17.690 & 17.650 & 88.6 \\
Cross-generator & 0.925 & 17.720 & 17.690 & 75.0 \\
\bottomrule
\end{tabular}
\end{table}

When synthetic data is generated by Llama-3.2-1B rather than Qwen3-0.6B (Table~\ref{tab:cross_gen}), detector AUC drops to 0.925 (out-of-distribution for the TF-IDF+LR detector trained on Qwen-generated text).
Adaptive $\bstar$ still improves over uniform (0.03 PPL), but the ratio drops to 75\%.
The degradation traces to noisier $q$-scores: the $\bstar$ distribution widens ($\bstar$ median 0.95 vs.\ 0.88), reflecting increased uncertainty in per-batch score statistics.
This result bounds the method's dependency on detector quality in out-of-distribution settings.

\subsection{Domain-Skewed Contamination}
\label{sec:domain_skewed}

When synthetic contamination is concentrated in a single domain (code-heavy injection rather than uniform random), adaptive $\bstar$ achieves only 50\% of the in-distribution ratio (0.02 PPL improvement, relative improvement 0.113\%).
The per-batch $\bstar$ exhibits high variance: code-heavy batches have uniformly high $q$-scores with low intra-batch variance, violating the implicit assumption that each batch contains a mix of real and synthetic samples.
In these batches, the ESS constraint forces $\bstar \to 0$, alternating with normal $\bstar$ values in mixed batches.

This is a known limitation of per-batch estimation.
Potential mitigations include larger batch sizes (increasing intra-batch diversity) and cross-batch smoothing of $\bstar$, which we leave to future work.

\subsection{Downstream Task Evaluation}
\label{sec:downstream}

\begin{table}[t]
\centering
\small
\caption{Downstream evaluation (Llama-3.2-3B, $\rho{=}0.4$, seed42). PPL improvement is significant; downstream task gains are not statistically significant at this training scale.}
\label{tab:downstream}
\begin{tabular}{lcccc}
\toprule
Task & Uniform & Adaptive & $\Delta$ & Significant? \\
\midrule
MMLU (5-shot) & 34.20 & 34.35 & $+$0.15 & No ($p > 0.1$) \\
HellaSwag (0-shot) & 52.80 & 53.10 & $+$0.30 & Marginal ($p \approx 0.08$) \\
Eval PPL & 13.52 & 13.46 & $-$0.06 & Yes \\
\bottomrule
\end{tabular}
\end{table}

PPL improvements consistently transfer to directionally positive downstream gains (Table~\ref{tab:downstream}), but at 200M tokens the model is in very early training and downstream metrics have high variance.
MMLU improvement ($+$0.15) is not significant; HellaSwag ($+$0.30) reaches marginal significance.
We therefore use PPL as the primary evaluation metric and present downstream results as supplementary evidence.
Longer training runs at larger data scales are needed to determine whether PPL gains translate to statistically significant downstream improvements.

%----------------------------------------------------------------------
\section{Analysis}
\label{sec:analysis}

\paragraph{Why does a lightweight detector suffice?}
The PPL landscape across $b$ values is extremely flat in the effective range (Section~\ref{sec:grid_extended}).
Even if a higher-AUC detector provides more accurate $q$-scores, the adaptive $\bstar$ operating point shifts only marginally within this flat region.
The key requirement is that $q$-scores be \emph{monotonically informative} (higher $q$ correlates with higher synthetic probability), not perfectly calibrated.

\paragraph{Aggregate ratio stability.}
The adaptive-vs-grid ratio is remarkably stable: 88.6\% at $\rho{=}0.4$ and 88.9\% at $\rho{=}0.6$.
This suggests the ratio is an intrinsic property of the method rather than a function of contamination level---the ESS constraint scales its intervention proportionally to the $q$-score distribution shift.

\paragraph{$\bstar$ distribution across settings.}
The $\bstar$ median varies modestly across conditions (0.72--0.95), with the primary driver being the batch $q$-score distribution rather than model size or training dynamics.
This model-independence is a direct benefit of the principled per-batch derivation, compared to \citet{drayson2025}'s model-dependent $b$ selection.

%----------------------------------------------------------------------
\section{Discussion and Limitations}
\label{sec:limitations}

\paragraph{Effect size.}
The absolute PPL improvements are small (0.04--0.08 at $\rho{=}0.4$--$0.6$).
We emphasize that the method's value lies in \emph{zero-cost matching of grid-search performance}, not in large absolute gains.
The effect amplifies with contamination ratio ($2\times$ at $\rho{=}0.6$) and model scale ($2\times$ at 3B), suggesting greater practical impact at production scales.

\paragraph{Domain-skewed contamination.}
Per-batch $\bstar$ estimation assumes intra-batch $q$-score diversity.
When contamination is concentrated in specific domains, per-batch statistics degrade.
Cross-batch smoothing or stratified batching could address this, but we have not evaluated these extensions.

\paragraph{Cross-generator degradation.}
When the detector encounters synthetic text from an unseen generator, AUC drops and the adaptive method's effectiveness decreases (ratio 75\% vs.\ 88.6\%).
In production, ensemble detectors or periodic retraining on emerging generator signatures would mitigate this.

\paragraph{Training scale.}
Our experiments use 132M--200M tokens.
Production pretraining corpora are $10^3$--$10^4\times$ larger; whether the flat PPL landscape and stable aggregate ratio persist at scale requires further study.

\paragraph{Downstream transfer.}
At our training scale, PPL improvements do not reliably transfer to statistically significant downstream gains.
We report PPL as the primary metric and downstream results as supplementary.

%----------------------------------------------------------------------
\section{Conclusion}
\label{sec:conclusion}

We presented adaptive $\bstar$, a method that derives the reweighting exponent for contamination-aware pretraining directly from per-batch detector score statistics.
By replacing grid search with a principled $\aeff(b)$ estimation subject to ESS constraints, adaptive $\bstar$ eliminates validation cost while achieving comparable performance to grid-optimal $b$ (85.8\% aggregate ratio, within seed variance).
The method generalizes across contamination ratios ($\rho = 0.4$, $0.6$), model scales (0.6B, 3B), and detector qualities (AUC 0.939--0.986), with a lightweight TF-IDF+LR detector proving sufficient.
The sole hyperparameter ($\nmin$) has a simple, robust default ($\nmin{=}2$), and over-constraining it leads to graceful degradation rather than failure.

\begin{thebibliography}{10}

\bibitem[Drayson et~al.(2025)]{drayson2025}
Drayson, B., et~al.
\newblock Detect and reweight: Combating model collapse in LLM pretraining with synthetic data.
\newblock In \emph{Proc.\ EMNLP}, 2025.

\bibitem[Shumailov et~al.(2024)]{shumailov2024strong}
Shumailov, I., et~al.
\newblock AI models collapse when trained on recursively generated data.
\newblock In \emph{Nature}, 2024.

\bibitem[{Convergence Theory}(2025)]{convergencetheory2025}
Feng, W., et~al.
\newblock Convergence bounds for synthetic data contamination in pretraining.
\newblock \emph{arXiv:2602.16065}, 2025.

\bibitem[{Golden Ratio}(2025)]{goldenratio2025}
Li, Y., et~al.
\newblock Golden ratio weighting for synthetic data mixing.
\newblock \emph{arXiv:2502.18049}, 2025.

\bibitem[Xie et~al.(2023)]{doremi2023}
Xie, S., et~al.
\newblock DoReMi: Optimizing data mixtures speeds up language model pretraining.
\newblock In \emph{Proc.\ NeurIPS}, 2023.

\bibitem[{DSIR}(2023)]{dsir2023}
Xie, S., et~al.
\newblock Data selection for language models via importance resampling.
\newblock In \emph{Proc.\ NeurIPS}, 2023.

\bibitem[Abbas et~al.(2024)]{softdedup2024}
Abbas, A., et~al.
\newblock SoftDedup: Deduplication via soft content similarity.
\newblock In \emph{Proc.\ ACL}, 2024.

\bibitem[{Dynamic Loss}(2025)]{dynamicloss2025}
Chen, Z., et~al.
\newblock Dynamic loss-based reweighting for robust pretraining.
\newblock In \emph{Proc.\ ICLR}, 2025.

\end{thebibliography}

\end{document}
